\section*{Решения}
\addcontentsline{toc}{section}{Решения}

далее матрица E будет составлена из векторов записанных в строку

\begin{enumerate}
    \item \textit{(Демо 2026, №1)} Составить систему линейных уравнений, задающую линейную оболочку векторов 
    \[
    a_1 = (1, 1, 2, 1, 2)^T, a_2 = (0, -1, -2, 1, -1)^T, a_3 = (3, 1, 2, 5, 4)^T
    \]    
    в $\mathbb{R}^5$

	построим матрицу где строки - данные векторы

	\[
		E = 
		\begin{pmatrix}
			1 & 1 & 2 & 1 & 2 \\
			0 & -1 & -2 & 1 & -1 \\
			3 & 1 & 2 & 5 & 4
		\end{pmatrix}
	\]

	ищем решения $Ex = 0$

	\[
		\begin{pmatrix}
			1 & 1 & 2 & 1 & 2 \\
			0 & -1 & -2 & 1 & -1 \\
			3 & 1 & 2 & 5 & 4
		\end{pmatrix}
		\times
		\begin{pmatrix}
			x_1 \\
			x_2 \\
			x_3 \\
			x_4 \\
			x_5
		\end{pmatrix}
		=0
	\]

	находим ФСР

	\[
		\begin{pmatrix}
			1 & 1 & 2 & 1 & 2 \\
			0 & -1 & -2 & 1 & -1 \\
			3 & 1 & 2 & 5 & 4
		\end{pmatrix}
		\underbrace{\Rightarrow}_{\text{ (3)-3(1) }} \\
		\begin{pmatrix}
			1 & 1 & 2 & 1 & 2 \\
			0 & -1 & -2 & 1 & -1 \\
			0 & -2 & -4 & 2 & -2
		\end{pmatrix}
		\underbrace{\Rightarrow}_{\text{ (3)-2(2) }} \\
		\begin{pmatrix}
			1 & 1 & 2 & 1 & 2 \\
			0 & -1 & -2 & 1 & -1 \\
			0 & 0 & 0 & 0 & 0
		\end{pmatrix}
	\]

	выражаем

	\[
		\begin{cases}
			x_1 + x_2 + 2 x_3 + x_4 + 2 x_5 = 0 \\
			-x_2 - 2 x_3 + x_4 - x_5 = 0
		\end{cases}
	\]

	выражаем главные:

	$x_2 = -2 x_3 + x_4 - x_5$

	$x_1 = -x_2 - 2 x_3 - x_4 - 2 x_5 = - 2 x_4 - x_5$

	уравнений будет 3:

	$x_3 = 0, x_4 = 0, x_5 = 1: x_2 = -1, x_1 = -1$

	$x_3 = 0, x_4 = 1, x_5 = 0: x_2 = 1, x_1 = -2$

	$x_3 = 1, x_4 = 0, x_5 = 0: x_2 = -2, x_1 = 0$

	ответ:

	\[
		\begin{cases}
			- x_1 - x_2 + x_5 = 0 \\
			-2 x_1 + x_2 + x_4 = 0 \\
			- 2 x_2 + x_3 = 0
		\end{cases}
	\]

    \item \textit{(Проскуряков, №1312)} Найти систему линейных уравнений, задающую линейное подпространство, натянутое на систему векторов:
    \[
    a_1 = (1, -1, 1, 0), \quad a_2 = (1, 1, 0, 1), \quad a_3 = (2, 0, 1, 1)
    \]

	\[
		E = 
		\begin{pmatrix}
			1 & -1 & 1 & 0 \\
			1 & 1 & 0 & 1 \\
			2 & 0 & 1 & 1 
		\end{pmatrix}
	\]

	находим ФСР

	\[
		\begin{pmatrix}
			1 & -1 & 1 & 0 \\
			1 & 1 & 0 & 1 \\
			2 & 0 & 1 & 1 
		\end{pmatrix}
		\underbrace{\Rightarrow}_{\text{ (2) - (1) }} \\
		\begin{pmatrix}
			1 & -1 & 1 & 0 \\
			0 & 2 & -1 & 1 \\
			2 & 0 & 1 & 1 
		\end{pmatrix}
		\underbrace{\Rightarrow}_{\text{ (3)-2(1) }} \\
		\begin{pmatrix}
			1 & -1 & 1 & 0 \\
			0 & 2 & -1 & 1 \\
			0 & 2 & -1 & 1 
		\end{pmatrix}
		\underbrace{\Rightarrow}_{\text{ (3)-(2) }} \\
		\begin{pmatrix}
			1 & -1 & 1 & 0 \\
			0 & 2 & -1 & 1 \\
			0 & 0 & 0 & 0 
		\end{pmatrix}
	\]

	$2 x_2 - x_3 + x_4 = 0$
	
	$x_1 - x_2 + x_3 = 0$

	$x_2 = \frac{x_3 - x_4}{2}$

	$x_1 = x_2 - x_3 = \frac{x_3 - x_4}{2} - x_3 = \frac{-x_3 - x_4}{2}$

	$x_3 = 0, x_4 = 1: x_2 = - \frac{1}{2}, x_1 = - \frac{1}{2}$

	$x_3 = 1, x_4 = 0: x_2 = \frac{1}{2}, x_1 = - \frac{1}{2}$

	ответ:

	\[
		\begin{cases}
			- \frac12 x_1 - \frac12 x_2 + x_4 = 0 \\
			- \frac12 x_1 + \frac12 x_2 + x_3 = 0
		\end{cases}
	\]

    \item \textit{(Проскуряков, №1313)} Найти систему линейных уравнений, задающую линейное подпространство, натянутое на систему векторов:
    \[
    a_1 = (1, -1, 1, -1, 1), \quad a_2 = (1, 1, 0, 0, 3), \quad a_3 = (3, 1, 1, -1, 7), \quad a_4 = (0, 2, -1, 1, 2)
    \]

	\[
		E = 
		\begin{pmatrix}
			1 & -1 & 1 & -1 & 1 \\
			1 & 1 & 0 & 0 & 3 \\
			3 & 1 & 1 & -1 & 7 \\
			0 & 2 & -1 & 1 & 2
		\end{pmatrix}
	\]

	находим ФСР

	\[
		\begin{pmatrix}
			1 & -1 & 1 & -1 & 1 \\
			1 & 1 & 0 & 0 & 3 \\
			3 & 1 & 1 & -1 & 7 \\
			0 & 2 & -1 & 1 & 2
		\end{pmatrix}
		\underbrace{\Rightarrow}_{\text{ (2) - (1) and 3 - 3(1)}} \\
		\begin{pmatrix}
			1 & -1 & 1 & -1 & 1 \\
			0 & 2 & -1 & 1 & 2 \\
			0 & 4 & -2 & 2 & 4 \\
			0 & 2 & -1 & 1 & 2
		\end{pmatrix}
	\]

	\[
		\begin{pmatrix}
			1 & -1 & 1 & -1 & 1 \\
			0 & 2 & -1 & 1 & 2 \\
			0 & 4 & -2 & 2 & 4 \\
			0 & 2 & -1 & 1 & 2
		\end{pmatrix}
		\underbrace{\Rightarrow}_{\text{ (4) - (2) and (3) - 2(2)}} \\
		\begin{pmatrix}
			1 & -1 & 1 & -1 & 1 \\
			0 & 2 & -1 & 1 & 2 \\
			0 & 0 & 0 & 0 & 0 \\
			0 & 0 & 0 & 0 & 0
		\end{pmatrix}
	\]

	$2 x_2 - x_3 + x_4  + 2 x_5 = 0$

	$x_1 - x_2 + x_3 - x_4 + x_5 = 0$

	$x_2 = \frac{x_3 - x_4 - 2 x_5}2$

	$x_1 = x_2 - x_3 + x_4 - x_5 = \frac{x_3 - x_4 - 2 x_5}2 - x_3 + x_4 - x_5 = \frac{-x_3 + x_4 - 4 x_5 }2$

	$x_3 = 0, x_4 = 0, x_5 = 1: x_2 = -1, x_1 = -2$

	$x_3 = 0, x_4 = 1, x_5 = 0: x_2 = - \frac12, x_1 = \frac12$

	$x_3 = 1, x_4 = 0, x_5 = 0: x_2 = \frac12, x_1 = - \frac12$

	ответ:

	\[
		\begin{cases}
			-2 x_1 - x_2 + x_5 = 0 \\
			x_1 - x_2 + 2 x_4 = 0 \\
			- x_1 + x_2 + 2 x_3 = 0
		\end{cases}
	\]


	\item \textit{(Кострикин, №35.11-а)} Найти базис и размерность линейной оболочки системы векторов в $\mathbb{R}^4$:
    \[
    a_1 = (1, 0, 0, -1), \quad a_2 = (2, 1, 1, 0), \quad a_3 = (1, 1, 1, 1), \quad a_4 = (1, 2, 3, 4), \quad a_5 = (0, 1, 2, 3)
    \]

    построим матрицу, где строки - данные векторы

    \[
        E =
        \begin{pmatrix}
            1 & 0 & 0 & -1 \\
            2 & 1 & 1 & 0 \\
            1 & 1 & 1 & 1 \\
            1 & 2 & 3 & 4 \\
            0 & 1 & 2 & 3
        \end{pmatrix}
    \]

    приводим к ступенчатому виду

    \[
        \begin{pmatrix}
            1 & 0 & 0 & -1 \\
            2 & 1 & 1 & 0 \\
            1 & 1 & 1 & 1 \\
            1 & 2 & 3 & 4 \\
            0 & 1 & 2 & 3
        \end{pmatrix}
        \underbrace{\Rightarrow}_{\text{ (2)-2(1), (3)-(1), (4)-(1) }} \\
        \begin{pmatrix}
            1 & 0 & 0 & -1 \\
            0 & 1 & 1 & 2 \\
            0 & 1 & 1 & 2 \\
            0 & 2 & 3 & 5 \\
            0 & 1 & 2 & 3
        \end{pmatrix}
    \]

	\[
        \begin{pmatrix}
            1 & 0 & 0 & -1 \\
            0 & 1 & 1 & 2 \\
            0 & 1 & 1 & 2 \\
            0 & 2 & 3 & 5 \\
            0 & 1 & 2 & 3
        \end{pmatrix}
        \underbrace{\Rightarrow}_{\text{ (3)-(2), (4)-2(2), (5)-(2) }} \\
        \begin{pmatrix}
            1 & 0 & 0 & -1 \\
            0 & 1 & 1 & 2 \\
            0 & 0 & 0 & 0 \\
            0 & 0 & 1 & 1 \\
            0 & 0 & 1 & 1
        \end{pmatrix}
	\]

	\[
		\begin{pmatrix}
		1 & 0 & 0 & -1 \\
		0 & 1 & 1 & 2 \\
		0 & 0 & 0 & 0 \\
		0 & 0 & 1 & 1 \\
		0 & 0 & 1 & 1
	\end{pmatrix}
	\underbrace{\Rightarrow}_{\text{ (5)-(4) и перестановка }} \\
	\begin{pmatrix}
		1 & 0 & 0 & -1 \\
		0 & 1 & 1 & 2 \\
		0 & 0 & 1 & 1 \\
		0 & 0 & 0 & 0 \\
		0 & 0 & 0 & 0
	\end{pmatrix}
	\]

    осталось 3 ненулевые строки, значит, ранг матрицы равен 3. размерность $\dim L = 3$.

    базис образуют векторы, соответствующие независимым строкам: первая, вторая и четвертая строки исходной матрицы ($a_3 = a_2 - a_1$ и $a_5 = a_4 - a_3$).

    ответ: размерность 3, базис: $a_1, a_2, a_4$
	тажке в качестве базиса можно было выписать ненулевые строчки полученные в Гауссе


    \item \textit{(Кострикин, №35.11-б)} Найти базис и размерность линейной оболочки системы векторов в $\mathbb{R}^5$:
    \[
    a_1 = (1, 1, 1, 1, 0), \quad a_2 = (1, 1, -1, -1, -1), \quad a_3 = (2, 2, 0, 0, -1),
    \]
    \[
    a_4 = (1, 1, 5, 5, 2), \quad a_5 = (1, -1, -1, 0, 0)
    \]

    построим матрицу

    \[
        E =
        \begin{pmatrix}
            1 & 1 & 1 & 1 & 0 \\
            1 & 1 & -1 & -1 & -1 \\
            2 & 2 & 0 & 0 & -1 \\
            1 & 1 & 5 & 5 & 2 \\
            1 & -1 & -1 & 0 & 0
        \end{pmatrix}
    \]

    приводим к ступенчатому виду

    \[
        \begin{pmatrix}
            1 & 1 & 1 & 1 & 0 \\
            1 & 1 & -1 & -1 & -1 \\
            2 & 2 & 0 & 0 & -1 \\
            1 & 1 & 5 & 5 & 2 \\
            1 & -1 & -1 & 0 & 0
        \end{pmatrix}
        \underbrace{\Rightarrow}_{\text{ (2)-(1), (3)-2(1), (4)-(1), (5)-(1) }} \\
        \begin{pmatrix}
            1 & 1 & 1 & 1 & 0 \\
            0 & 0 & -2 & -2 & -1 \\
            0 & 0 & -2 & -2 & -1 \\
            0 & 0 & 4 & 4 & 2 \\
            0 & -2 & -2 & -1 & 0
        \end{pmatrix}
    \]

	\[
        \begin{pmatrix}
            1 & 1 & 1 & 1 & 0 \\
            0 & 0 & -2 & -2 & -1 \\
            0 & 0 & -2 & -2 & -1 \\
            0 & 0 & 4 & 4 & 2 \\
            0 & -2 & -2 & -1 & 0
        \end{pmatrix}
        \underbrace{\Rightarrow}_{\text{ (3)-(2), (4)+2(2) }} \\
        \begin{pmatrix}
            1 & 1 & 1 & 1 & 0 \\
            0 & 0 & -2 & -2 & -1 \\
            0 & 0 & 0 & 0 & 0 \\
            0 & 0 & 0 & 0 & 0 \\
            0 & -2 & -2 & -1 & 0
        \end{pmatrix}
	\]


	\[
        \begin{pmatrix}
            1 & 1 & 1 & 1 & 0 \\
            0 & 0 & -2 & -2 & -1 \\
            0 & 0 & 0 & 0 & 0 \\
            0 & 0 & 0 & 0 & 0 \\
            0 & -2 & -2 & -1 & 0
        \end{pmatrix}
        \underbrace{\Rightarrow}_{\text{ меняем (2) и (5) местами }} \\
        \begin{pmatrix}
            1 & 1 & 1 & 1 & 0 \\
            0 & -2 & -2 & -1 & 0 \\
            0 & 0 & -2 & -2 & -1 \\
            0 & 0 & 0 & 0 & 0 \\
            0 & 0 & 0 & 0 & 0
        \end{pmatrix}
	\]

    осталось 3 ненулевые строки, ранг матрицы равен 3. размерность $\dim L = 3$.

    базис образуют векторы: первая строка ($a_1$), переставленная пятая строка ($a_5$) и вторая строка ($a_2$). остальные зависят от них ($a_3 = a_1 + a_2$, $a_4 = 3a_1 - 2a_2$).

    ответ: размерность 3, базис: $a_1, a_2, a_5$

	тажке в качестве базиса можно было выписать ненулевые строчки полученные в Гауссе


	\item \textit{(Демо 2026, №2)} Найти размерности и базисы суммы и пересечения подпространств $V_1, V_2$ в $\mathbb{R}^4$,
    \[
    V_1 = \langle a_1, a_2, a_3 \rangle,\qquad
    V_2 = \langle b_1, b_2 \rangle,
    \]
    где
    \[
    a_1 = (1, 0, -3, -2)^T,\quad
    a_2 = (7, 1, 9, 14)^T,\quad
    a_3 = (-4, 1, 2, -9)^T,
    \]
    \[
    b_1 = (10, 1, 0, 8)^T,\quad
    b_2 = (-3, 0, 1, -3)^T.
    \]

    найдём размерности и базисы $V_1$ и $V_2$

    для $V_1$ запишем векторы по строкам:
    \[
        \begin{pmatrix}
            1 & 0 & -3 & -2 \\
            7 & 1 & 9 & 14 \\
            -4 & 1 & 2 & -9
        \end{pmatrix}
        \underbrace{\Rightarrow}_{\text{ (2) - 7(1) и (3) + 4(1) }} \\
        \begin{pmatrix}
            1 & 0 & -3 & -2 \\
            0 & 1 & 30 & 28 \\
            0 & 1 & -10 & -17
        \end{pmatrix}
        \underbrace{\Rightarrow}_{\text{ (3) - (2) }} \\
        \begin{pmatrix}
            1 & 0 & -3 & -2 \\
            0 & 1 & 30 & 28 \\
            0 & 0 & -40 & -45
        \end{pmatrix}
    \]
    $\dim V_1 = 3$, базис $V_1$: $a_1, a_2, a_3$ (исходные векторы линейно независимы).

    для $V_2$:
    \[
        \begin{pmatrix}
            10 & 1 & 0 & 8 \\
            -3 & 0 & 1 & -3
        \end{pmatrix}
        \underbrace{\Rightarrow}_{\text{ (1) + 3(2) }} \\
        \begin{pmatrix}
            1 & 1 & 3 & -1 \\
            -3 & 0 & 1 & -3
        \end{pmatrix}
    \]
    строки не пропорциональны, $\dim V_2 = 2$, базис $V_2$: $b_1, b_2$.

    найдём сумму $V_1 + V_2$:

    запишем в матрицу базисы $V_1$ (возьмем упрощенный ступенчатый из первого пункта для удобства) и векторы $V_2$:
    \[
        \begin{pmatrix}
            1 & 0 & -3 & -2 \\
            0 & 1 & 30 & 28 \\
            0 & 0 & -40 & -45 \\
            10 & 1 & 0 & 8 \\
            -3 & 0 & 1 & -3
        \end{pmatrix}
        \underbrace{\Rightarrow}_{\text{ (4) - 10(1) и (5) + 3(1) }} \\
        \begin{pmatrix}
            1 & 0 & -3 & -2 \\
            0 & 1 & 30 & 28 \\
            0 & 0 & -40 & -45 \\
            0 & 1 & 30 & 28 \\
            0 & 0 & -8 & -9
        \end{pmatrix}
        \underbrace{\Rightarrow}_{\text{ (4) - (2) и (5) - 1/5(3) }} \\
        \begin{pmatrix}
            1 & 0 & -3 & -2 \\
            0 & 1 & 30 & 28 \\
            0 & 0 & -40 & -45 \\
            0 & 0 & 0 & 0 \\
            0 & 0 & 0 & 0
        \end{pmatrix}
    \]
    векторы $b_1$ и $b_2$ полностью обнулились. $\dim(V_1 + V_2) = 3$. 
    вывод: $V_2 \subset V_1$, поэтому сумма совпадает с $V_1$, базис суммы: $a_1, a_2, a_3$.

	здесь можно было явно заметить что одно подпространство входит в другое и пересечение очевидно

	однако для подготовки к экзамену приведены два способа поиска базиса пересечение в общем случае:

    найдём пересечение $V_1 \cap V_2$ (Способ 1: через линейные комбинации)

    вектор $x \in V_1 \cap V_2$ выражается через оба базиса:
    $\alpha_1 a_1 + \alpha_2 a_2 + \alpha_3 a_3 = \beta_1 b_1 + \beta_2 b_2$
    $\alpha_1 a_1 + \alpha_2 a_2 + \alpha_3 a_3 - \beta_1 b_1 - \beta_2 b_2 = 0$

    запишем векторы \textbf{по столбцам} и найдем ФСР для коэффициентов:
    \[
        \begin{pmatrix}
            1 & 7 & -4 & -10 & 3 \\
            0 & 1 & 1 & -1 & 0 \\
            -3 & 9 & 2 & 0 & -1 \\
            -2 & 14 & -9 & -8 & 3
        \end{pmatrix}
        \underbrace{\Rightarrow}_{\text{ (3)+3(1), (4)+2(1) }} \\
        \begin{pmatrix}
            1 & 7 & -4 & -10 & 3 \\
            0 & 1 & 1 & -1 & 0 \\
            0 & 30 & -10 & -30 & 8 \\
            0 & 28 & -17 & -28 & 9
        \end{pmatrix}
    \]

	\[
        \begin{pmatrix}
            1 & 7 & -4 & -10 & 3 \\
            0 & 1 & 1 & -1 & 0 \\
            0 & 30 & -10 & -30 & 8 \\
            0 & 28 & -17 & -28 & 9
        \end{pmatrix}
        \underbrace{\Rightarrow}_{\text{ (3)-30(2), (4)-28(2) }} \\
        \begin{pmatrix}
            1 & 7 & -4 & -10 & 3 \\
            0 & 1 & 1 & -1 & 0 \\
            0 & 0 & -40 & 0 & 8 \\
            0 & 0 & -45 & 0 & 9
        \end{pmatrix}
	\]

	\[
        \begin{pmatrix}
            1 & 7 & -4 & -10 & 3 \\
            0 & 1 & 1 & -1 & 0 \\
            0 & 0 & -40 & 0 & 8 \\
            0 & 0 & -45 & 0 & 9
        \end{pmatrix}
        \underbrace{\Rightarrow}_{\text{ делим (3) на -8, (4) на -9 }} \\
        \begin{pmatrix}
            1 & 7 & -4 & -10 & 3 \\
            0 & 1 & 1 & -1 & 0 \\
            0 & 0 & 5 & 0 & -1 \\
            0 & 0 & 5 & 0 & -1
        \end{pmatrix}
	\]

    четвертая строка обнуляется. свободные переменные: $\beta_1$ и $\beta_2$.
    так как на $\beta_1$ и $\beta_2$ нет никаких ограничений (они могут быть любыми), значит \textbf{любой} вектор из $V_2$ принадлежит $V_1$. 
    $\dim(V_1 \cap V_2) = 2$, базис: $b_1, b_2$.


    найдём пересечение $V_1 \cap V_2$ (Способ 2: через системы уравнений)
    
    найдем уравнение $V_1$. запишем $a_1, a_2, a_3$ по столбцам и припишем столбец $(x_1, x_2, x_3, x_4)^T$:
    \[
        \begin{pmatrix}
            1 & 7 & -4 & | & x_1 \\
            0 & 1 & 1 & | & x_2 \\
            -3 & 9 & 2 & | & x_3 \\
            -2 & 14 & -9 & | & x_4
        \end{pmatrix}
        \Rightarrow \dots \Rightarrow
        \begin{pmatrix}
            1 & 7 & -4 & | & x_1 \\
            0 & 1 & 1 & | & x_2 \\
            0 & 0 & -40 & | & 3x_1 - 30x_2 + x_3 \\
            0 & 0 & 0 & | & 11x_1 - 46x_2 + 9x_3 - 8x_4
        \end{pmatrix}
    \]
    уравнение $V_1$: $11x_1 - 46x_2 + 9x_3 - 8x_4 = 0$

    аналогично найдем уравнения $V_2$ (векторы $b_1, b_2$ по столбцам):
    \[
        \begin{pmatrix}
            10 & -3 & | & x_1 \\
            1 & 0 & | & x_2 \\
            0 & 1 & | & x_3 \\
            8 & -3 & | & x_4
        \end{pmatrix}
        \Rightarrow \dots \Rightarrow
        \begin{pmatrix}
            1 & 0 & | & x_2 \\
            0 & 1 & | & x_3 \\
            0 & 0 & | & x_1 - 10x_2 + 3x_3 \\
            0 & 0 & | & -8x_2 + 3x_3 + x_4
        \end{pmatrix}
    \]
    система $V_2$: 
    $\begin{cases} x_1 - 10x_2 + 3x_3 = 0 \\ -8x_2 + 3x_3 + x_4 = 0 \end{cases}$

    объединяем в общую систему $V_1 \cap V_2$:
    \[
        \begin{cases} 
            11x_1 - 46x_2 + 9x_3 - 8x_4 = 0 \\ 
            x_1 - 10x_2 + 3x_3 = 0 \\ 
            -8x_2 + 3x_3 + x_4 = 0 
        \end{cases}
    \]
    заметим, что первое уравнение является линейной комбинацией второго и третьего: $11(1) - 8(2)$. оно обнуляется при Гауссе. 
    система сводится к уравнениям $V_2$, значит $\dim(V_1 \cap V_2) = 4 - 2 = 2$.
    ФСР этой системы даст исходные векторы $b_1, b_2$.

    \textbf{Ответ:} 
    $\dim(V_1 + V_2) = 3$, базис: $a_1, a_2, a_3$
    $\dim(V_1 \cap V_2) = 2$, базис: $b_1, b_2$


    \item \textit{(Проскуряков, №1317)} Найти размерность $s$ суммы и размерность $d$ пересечения линейных подпространств $L_1 = \langle a_1, a_2 \rangle$ и $L_2 = \langle b_1, b_2 \rangle$ в $\mathbb{R}^4$, где
    \[
    a_1 = (1, 2, 0, 1)^T, \quad a_2 = (1, 1, 1, 0)^T,
    \]
    \[
    b_1 = (1, 0, 1, 0)^T, \quad b_2 = (1, 3, 0, 1)^T.
    \]

	найдём $\dim L_1$:

	\[
		\begin{pmatrix}
			1 & 2 & 0 & 1 \\
			1 & 1 & 1 & 0
		\end{pmatrix}
        \underbrace{\Rightarrow}_{\text{ (2) - (1) }} \\
		\begin{pmatrix}
			1 & 2 & 0 & 1 \\
			0 & -1 & 1 & -1
		\end{pmatrix}
	\]

	$\dim L_1 = 2$

	найдём $\dim L_2$:

	\[
		\begin{pmatrix}
			1 & 0 & 1 & 0 \\
			1 & 3 & 0 & 1
		\end{pmatrix}
        \underbrace{\Rightarrow}_{\text{ (2) - (1) }} \\
		\begin{pmatrix}
			1 & 0 & 1 & 0 \\
			0 & 3 & -1 & 1
		\end{pmatrix}
	\]

	$\dim L_2 = 2$

	$\dim (L_1 + L_2)$:

	\[
		\begin{pmatrix}
			1 & 2 & 0 & 1 \\
			0 & -1 & 1 & -1 \\
			1 & 0 & 1 & 0 \\
			0 & 3 & -1 & 1
		\end{pmatrix}
        \underbrace{\Rightarrow}_{\text{ (3) - (1) and (4) + 3(2)}} \\
		\begin{pmatrix}
			1 & 2 & 0 & 1 \\
			0 & -1 & 1 & -1 \\
			0 & -2 & 1 & -1 \\
			0 & 0 & 2 & -2
		\end{pmatrix}
	\]

	\[
		\begin{pmatrix}
			1 & 2 & 0 & 1 \\
			0 & -1 & 1 & -1 \\
			0 & -2 & 1 & -1 \\
			0 & 0 & 2 & -2
		\end{pmatrix}
        \underbrace{\Rightarrow}_{\text{ (3) - 2(2)  }} \\
		\begin{pmatrix}
			1 & 2 & 0 & 1 \\
			0 & -1 & 1 & -1 \\
			0 & 0 & -1 & 1 \\
			0 & 0 & 2 & -2
		\end{pmatrix}
	\]

	\[
		\begin{pmatrix}
			1 & 2 & 0 & 1 \\
			0 & -1 & 1 & -1 \\
			0 & 0 & -1 & 1 \\
			0 & 0 & 2 & -2
		\end{pmatrix}
        \underbrace{\Rightarrow}_{\text{ (4) + 2(3)  }} \\
		\begin{pmatrix}
			1 & 2 & 0 & 1 \\
			0 & -1 & 1 & -1 \\
			0 & 0 & -1 & 1 \\
			0 & 0 & 0 & 0
		\end{pmatrix}
	\]

	$\dim (L_1 + L_2) = 3$

	тогда $\dim (L_1 \cap L_2) = 2 + 2 - 3 = 1$

	\item \textit{(Проскуряков, №1318)} Найти размерность $s$ суммы и размерность $d$ пересечения линейных подпространств $L_1 = \langle a_1, a_2, a_3 \rangle$ и $L_2 = \langle b_1, b_2, b_3 \rangle$ в $\mathbb{R}^4$, где
        \[
        a_1 = (1, 1, 1, 1)^T, \quad a_2 = (1, -1, 1, -1)^T, \quad a_3 = (1, 3, 1, 3)^T,
        \]
        \[
        b_1 = (1, 2, 0, 2)^T, \quad b_2 = (1, 2, 1, 2)^T, \quad b_3 = (3, 1, 3, 1)^T.
        \]

    найдём $\dim L_1$:
    \[
        \begin{pmatrix}
            1 & 1 & 1 & 1 \\
            1 & -1 & 1 & -1 \\
            1 & 3 & 1 & 3
        \end{pmatrix}
        \underbrace{\Rightarrow}_{\text{ (2)-(1), (3)-(1) }} \\
        \begin{pmatrix}
            1 & 1 & 1 & 1 \\
            0 & -2 & 0 & -2 \\
            0 & 2 & 0 & 2
        \end{pmatrix}
        \underbrace{\Rightarrow}_{\text{ (3)+(2) }} \\
        \begin{pmatrix}
            1 & 1 & 1 & 1 \\
            0 & -2 & 0 & -2 \\
            0 & 0 & 0 & 0
        \end{pmatrix}
    \]
    $\dim L_1 = 2$

    найдём $\dim L_2$:
    \[
        \begin{pmatrix}
            1 & 2 & 0 & 2 \\
            1 & 2 & 1 & 2 \\
            3 & 1 & 3 & 1
        \end{pmatrix}
        \underbrace{\Rightarrow}_{\text{ (2)-(1), (3)-3(1) }} \\
        \begin{pmatrix}
            1 & 2 & 0 & 2 \\
            0 & 0 & 1 & 0 \\
            0 & -5 & 3 & -5
        \end{pmatrix}
        \underbrace{\Rightarrow}_{\text{ (2) <-> (3) }} \\
        \begin{pmatrix}
            1 & 2 & 0 & 2 \\
            0 & -5 & 3 & -5 \\
            0 & 0 & 1 & 0
        \end{pmatrix}
    \]
    $\dim L_2 = 3$

    найдём $\dim(L_1 + L_2)$:
    \[
        \begin{pmatrix}
            1 & 1 & 1 & 1 \\
            1 & -1 & 1 & -1 \\
            1 & 2 & 0 & 2 \\
            1 & 2 & 1 & 2 \\
            3 & 1 & 3 & 1
        \end{pmatrix}
        \underbrace{\Rightarrow}_{\text{ вычитаем (1) из остальных }} \\
        \begin{pmatrix}
            1 & 1 & 1 & 1 \\
            0 & -2 & 0 & -2 \\
            0 & 1 & -1 & 1 \\
            0 & 1 & 0 & 1 \\
            0 & -2 & 0 & -2
        \end{pmatrix}
        \underbrace{\Rightarrow}_{\text{ (5)-(2) и делим (2) на -2 }} \\
        \begin{pmatrix}
            1 & 1 & 1 & 1 \\
            0 & 1 & 0 & 1 \\
            0 & 1 & -1 & 1 \\
            0 & 1 & 0 & 1 \\
            0 & 0 & 0 & 0
        \end{pmatrix}
        \underbrace{\Rightarrow}_{\text{ (3)-(2), (4)-(2) }} \\
        \begin{pmatrix}
            1 & 1 & 1 & 1 \\
            0 & 1 & 0 & 1 \\
            0 & 0 & -1 & 0 \\
            0 & 0 & 0 & 0 \\
            0 & 0 & 0 & 0
        \end{pmatrix}
    \]
    $\dim(L_1 + L_2) = 3$ ($s = 3$)

    по формуле Грассмана: 
    $\dim(L_1 \cap L_2) = \dim L_1 + \dim L_2 - \dim(L_1 + L_2) = 2 + 3 - 3 = 2$ ($d = 2$).

	\item \textit{(Демо 2026, №3)} Доказать, что пространство $\mathbb{R}^4$ является прямой суммой подпространств
    \[
    L_1=\langle a_1,a_2\rangle,\qquad L_2=\langle b_1,b_2\rangle,
    \]
    и разложить вектор $x=(0,-2,2,0)^T$ на сумму проекций на эти подпространства, где
    \[
    a_1=(1,1,1,0)^T,\quad a_2=(1,1,0,1)^T,\quad
    b_1=(1,0,1,1)^T,\quad b_2=(1,1,-1,-1)^T.
    \]

    доказательство прямой суммы

    чтобы $\mathbb{R}^4 = L_1 \oplus L_2$, необходимо и достаточно, чтобы $\dim(L_1 + L_2) = 4$.
    запишем образующие векторы по строкам и найдем ранг:

    \[
        \begin{pmatrix}
            1 & 1 & 1 & 0 \\
            1 & 1 & 0 & 1 \\
            1 & 0 & 1 & 1 \\
            1 & 1 & -1 & -1
        \end{pmatrix}
        \underbrace{\Rightarrow}_{\text{ (2)-(1), (3)-(1), (4)-(1) }} \\
        \begin{pmatrix}
            1 & 1 & 1 & 0 \\
            0 & 0 & -1 & 1 \\
            0 & -1 & 0 & 1 \\
            0 & 0 & -2 & -1
        \end{pmatrix}
    \]

    \[
        \begin{pmatrix}
            1 & 1 & 1 & 0 \\
            0 & 0 & -1 & 1 \\
            0 & -1 & 0 & 1 \\
            0 & 0 & -2 & -1
        \end{pmatrix}
        \underbrace{\Rightarrow}_{\text{ (2) <-> (3) }} \\
        \begin{pmatrix}
            1 & 1 & 1 & 0 \\
            0 & -1 & 0 & 1 \\
            0 & 0 & -1 & 1 \\
            0 & 0 & -2 & -1
        \end{pmatrix}
    \]

    \[
        \begin{pmatrix}
            1 & 1 & 1 & 0 \\
            0 & -1 & 0 & 1 \\
            0 & 0 & -1 & 1 \\
            0 & 0 & -2 & -1
        \end{pmatrix}
        \underbrace{\Rightarrow}_{\text{ (4) - 2(3) }} \\
        \begin{pmatrix}
            1 & 1 & 1 & 0 \\
            0 & -1 & 0 & 1 \\
            0 & 0 & -1 & 1 \\
            0 & 0 & 0 & -3
        \end{pmatrix}
    \]

    ранг матрицы равен 4, следовательно векторы образуют базис $\mathbb{R}^4$.
    значит, $\dim(L_1 + L_2) = 4$. так как $\dim L_1 \le 2$ и $\dim L_2 \le 2$, то их сумма прямая ($L_1 \cap L_2 = \{0\}$).

    разложение вектора $x$

    представим $x = x_1 + x_2$, где $x_1 \in L_1$ и $x_2 \in L_2$.
    тогда $x = \alpha_1 a_1 + \alpha_2 a_2 + \beta_1 b_1 + \beta_2 b_2$.
    составим неоднородную СЛАУ, записав векторы по столбцам:

    \[
        \begin{pmatrix}
            1 & 1 & 1 & 1 & | & 0 \\
            1 & 1 & 0 & 1 & | & -2 \\
            1 & 0 & 1 & -1 & | & 2 \\
            0 & 1 & 1 & -1 & | & 0
        \end{pmatrix}
        \underbrace{\Rightarrow}_{\text{ (2)-(1), (3)-(1) }} \\
        \begin{pmatrix}
            1 & 1 & 1 & 1 & | & 0 \\
            0 & 0 & -1 & 0 & | & -2 \\
            0 & -1 & 0 & -2 & | & 2 \\
            0 & 1 & 1 & -1 & | & 0
        \end{pmatrix}
    \]

    из 2-й строки: $-\beta_1 = -2 \Rightarrow \beta_1 = 2$.
    из 4-й строки: $\alpha_2 + \beta_1 - \beta_2 = 0 \Rightarrow \alpha_2 + 2 - \beta_2 = 0 \Rightarrow \alpha_2 - \beta_2 = -2$.
    из 3-й строки: $-\alpha_2 - 2\beta_2 = 2 \Rightarrow \alpha_2 + 2\beta_2 = -2$.

    вычтем уравнения: $(\alpha_2 - \beta_2) - (\alpha_2 + 2\beta_2) = -2 - (-2) \Rightarrow -3\beta_2 = 0 \Rightarrow \beta_2 = 0$.
    тогда $\alpha_2 - 0 = -2 \Rightarrow \alpha_2 = -2$.
    из 1-й строки: $\alpha_1 + \alpha_2 + \beta_1 + \beta_2 = 0 \Rightarrow \alpha_1 - 2 + 2 + 0 = 0 \Rightarrow \alpha_1 = 0$.

    находим проекции:
    $x_1 = \alpha_1 a_1 + \alpha_2 a_2 = 0 \cdot a_1 - 2 \cdot (1, 1, 0, 1)^T = (-2, -2, 0, -2)^T$
    $x_2 = \beta_1 b_1 + \beta_2 b_2 = 2 \cdot (1, 0, 1, 1)^T + 0 \cdot b_2 = (2, 0, 2, 2)^T$

    ответ: $x_1 = (-2, -2, 0, -2)^T$, $x_2 = (2, 0, 2, 2)^T$.

\end{enumerate}
